% Generated draft: requires mathematical and source review.
\documentclass[11pt]{article}
\usepackage[margin=25mm]{geometry}
\usepackage{fontspec}
\setmainfont{Latin Modern Roman}
\usepackage{amsmath,amssymb,mathtools,unicode-math}
\setmathfont{Latin Modern Math}
\usepackage{enumitem,needspace,xurl,hyperref}
\setlength{\parindent}{0pt}
\setlength{\parskip}{5pt}
\setlength{\emergencystretch}{3em}
\newcommand{\N}{\mathbb{N}}
\newcommand{\Z}{\mathbb{Z}}
\newcommand{\Q}{\mathbb{Q}}
\newcommand{\R}{\mathbb{R}}
\newcommand{\C}{\mathbb{C}}
\newcommand{\F}{\mathbb{F}}
\newcommand{\A}{\mathbb{A}}
\newcommand{\PP}{\mathbb{P}}
\newcommand{\E}{\mathbb{E}}
\newcommand{\eps}{\varepsilon}
\newcommand{\dd}{\,\mathrm{d}}
\begin{document}
\section*{UnsolvedMath ID: 30005223}
\textbf{Source problem code:} catalogue entry
\subsection{Definitions and mathematical statement}
For an integer \(n\geq 1\), let
\[
\mathcal P(n)=\{(\lambda_1,\ldots,\lambda_r):r\geq 1,\ \lambda_1\geq\cdots\geq\lambda_r\geq 1,\ \lambda_i\in\Z,\ \textstyle\sum_i\lambda_i=n\},
\qquad p(n)=|\mathcal P(n)|.
\]
Let \(S_n\) be the symmetric group on \(n\) letters. Its ordinary irreducible complex characters are indexed by \(\lambda\in\mathcal P(n)\); write \(\chi_\lambda\) for the corresponding character. For \(\mu\in\mathcal P(n)\), let \(K_\mu\subseteq S_n\) be the conjugacy class of permutations with cycle type \(\mu\), and define \(\chi_\lambda(\mu)=\chi_\lambda(g)\) for any \(g\in K_\mu\).

Under the interpretation that \(\lambda\) and \(\mu\) are independent and each is uniform on \(\mathcal P(n)\), define
\[
P_n=\Pr\bigl(\chi_\lambda(\mu)=0\bigr)
=\frac{\bigl|\{(\lambda,\mu)\in\mathcal P(n)^2:\chi_\lambda(\mu)=0\}\bigr|}{p(n)^2}.
\]
\textbf{Problem.} Determine the asymptotic behavior of \(P_n\) as \(n\to\infty\). In particular, does \(\lim_{n\to\infty}P_n\) exist, and, if so, what is its value? The supplied literature assessment specifies uniform partitions and conjugacy classes. Independence is the natural product-sampling interpretation used in this formulation, but is not expressly stated in the extracted question; this qualification requires source review. No proposed value of the limit is supplied.

\textbf{Distinction from the supplied background results.} Uniform choice of a conjugacy class is different from choosing an element uniformly from \(S_n\). With independent uniform choices of \(\chi\in\operatorname{Irr}(S_n)\) and \(g\in S_n\), the supplied excerpt reports that \(\Pr(\chi(g)=0)\to 1\). It also reports an analogous result for \(\operatorname{GL}(n,q)\): if \(P_{n,q}\) is the proportion of pairs \((\chi,g)\in\operatorname{Irr}(\operatorname{GL}(n,q))\times\operatorname{GL}(n,q)\) with \(\chi(g)=0\), then \(\inf_q P_{n,q}\to 1\) as \(n\to\infty\), where \(q\) ranges over prime powers. These are supplied contextual assertions, not independently verified results here, and do not answer the uniform-class question. The supplied \(S_4\) example gives the uniform-class probability \(4/25=0.16\); its element-sampling numerical expression is damaged or inconsistent and is left unresolved.
\subsection{Short English statement}
Choose two partitions \(\lambda,\mu\) of \(n\) independently and uniformly. Determine how often the irreducible character \(\chi_\lambda\) vanishes on the conjugacy class of cycle type \(\mu\) as \(n\) grows. Independence needs confirmation from the source. The supplied open-status assessment concerns uniform classes; the background result about uniform group elements uses a different probability measure.
\subsection{Sources}
\begin{itemize}
\item Local UnsolvedMath registry, supplied source record with permanent ID \(30005223\), title \emph{Limiting Vanishing Probability for Symmetric Group Characters}, exact source problem code \texttt{catalogue entry}. The registry records proposed year 2022 and supplies no proposer. Its labels \texttt{open} and \texttt{corrected\_verified} are supplied curation labels, not independent verification here. The accompanying verification note states that only formula delimiters, notation, grammar, and typography were repaired, with mathematical content and scope preserved.
\item Supplied report citation: \emph{Character Theory and Categorification (2023)}, Oberwolfach Reports, DOI 10.4171/owr/2022/39, pp.\ 2267--2333. \url{https://doi.org/10.4171/owr/2022/39}. The supplied background describes these reports as records of mathematical workshops and open-problem discussions.
\item Supplied literature source: Miller, \emph{Dense proportions of zeros}. \url{https://doi.org/10.1016/j.crma.2019.10.006}.
\item Supplied literature source: \emph{Oberwolfach Report 39/2022}. \url{https://ems.press/content/serial-article-files/46974}.
\item Literature assessment dated exactly \texttt{2026-08-22}: the supplied curation labels the problem open and says that joint character-and-element vanishing probabilities tend to one in several settings, while the limiting behavior for uniform partitions and conjugacy classes of symmetric groups is explicitly left open. This is supplied dated literature triage derived from the OWR curation, not an independently verified assessment; time-sensitive status and cited claims require independent checking.
\item The extracted background attributes a bound involving character degrees and conjugacy-class sizes to joint work of P.\ X.\ Gallagher, M.\ J.\ Larsen, and an unnamed speaker, cited there as reference [1]. It attributes the symmetric-group result and a class-covering inequality to reference [3]. The supplied record does not identify these numbered references explicitly enough to reproduce their bibliographic entries.
\end{itemize}
\end{document}
